% !Mode:: "TeX:UTF-8"
% !TEX root = ../../main.tex
% ============================================================
% TOPIC 3. MINORS, COFACTORS AND DETERMINANTS OF ORDER n
% ============================================================
\mashgulot{3}{Minors and cofactors. Evaluation of determinants of order \texorpdfstring{$n$}{n}}
\begin{tarif}
Let $a_{ij}$ be an entry of a determinant of order $n$. The determinant of order
$n-1$ obtained by deleting the $i$th row and the $j$th column, which contain
$a_{ij}$, is called the \textbf{minor} of $a_{ij}$ and is denoted by $M_{ij}$.
\end{tarif}
\begin{misol}
For the matrix
$A=\begin{pmatrix}2 & 1 & 3\\ 0 & -1 & 2\\ 4 & 3 & 1\end{pmatrix}$
we find the minor of the entry $a_{12}=1$. Deleting the first row and the second
column, we obtain
\[\kf{M_{12} = \begin{vmatrix} 0 & 2 \\ 4 & 1 \end{vmatrix} = 0\cdot 1 - 2\cdot 4 = -8.}\]
\end{misol}
\begin{misol}
In the same matrix the minor of $a_{23}=2$ is
\[\kf{M_{23} = \begin{vmatrix} 2 & 1 \\ 4 & 3 \end{vmatrix} = 2\cdot 3 - 1\cdot 4 = 2.}\]
\end{misol}
\begin{misol}
The minor of $a_{31}=4$ is
\[\kf{M_{31} = \begin{vmatrix} 1 & 3 \\ -1 & 2 \end{vmatrix} = 1\cdot 2 - 3\cdot(-1) = 2+3 = 5.}\]
\end{misol}
\textbf{Cofactors.}
\vspace{3pt}
\begin{tarif}
The \textbf{cofactor} (algebraic complement) of the entry $a_{ij}$ is the number
\[\kf{A_{ij} = (-1)^{i+j} \cdot M_{ij},}\]
where $M_{ij}$ is the minor of $a_{ij}$; the factor $(-1)^{i+j}$ determines the sign.
\end{tarif}

\begin{misol}
For the matrix $A$ above,
\begin{align*}
A_{12} &= (-1)^{1+2}\cdot M_{12} = (-1)^3 \cdot (-8) = 8, \\
A_{23} &= (-1)^{2+3}\cdot M_{23} = (-1)^5 \cdot 2 = -2, \\
A_{31} &= (-1)^{3+1}\cdot M_{31} = (-1)^4 \cdot 5 = 5.
\end{align*}
\end{misol}
\vspace{3pt}
\begin{teorema}[Laplace expansion]
A determinant of order $n$ equals the sum of the products of the entries of any
of its rows (or columns) by their cofactors. Expansion along the $i$th row reads
\[\kf{\det A = \sum\limits_{j=1}^{n} a_{ij}\cdot A_{ij}
= a_{i1}A_{i1} + a_{i2}A_{i2} + \cdots + a_{in}A_{in},}\]
and expansion along the $j$th column reads
\[\kf{\det A = \sum\limits_{i=1}^{n} a_{ij}\cdot A_{ij}
= a_{1j}A_{1j} + a_{2j}A_{2j} + \cdots + a_{nj}A_{nj}.}\]
\end{teorema}
The theorem shows that every row and every column give the same value, so in
practice we expand along the row or column containing \textbf{the most zeros}.
A closely related fact is often used: if we multiply the entries of one row by the
cofactors of a \emph{different} row and add, the result is zero,
\[\kf{\sum\limits_{j=1}^{n} a_{ij}\cdot A_{kj} = 0, \quad i\neq k.}\]
Indeed, this sum is the expansion along the $k$th row of the determinant in which
the $k$th row has been replaced by the $i$th one, and that determinant has two
equal rows.
\begin{misol}
We expand the determinant along the second column:
\[\kf{\begin{vmatrix}
2 & 1 & 3 \\
0 & -1 & 2 \\
4 & 3 & 1
\end{vmatrix}
= a_{12}A_{12} + a_{22}A_{22} + a_{32}A_{32}}\]
\[\kf{= 1\cdot(-1)^{1+2}\begin{vmatrix}0&2\\4&1\end{vmatrix}
+ (-1)\cdot(-1)^{2+2}\begin{vmatrix}2&3\\4&1\end{vmatrix}
+ 3\cdot(-1)^{3+2}\begin{vmatrix}2&3\\0&2\end{vmatrix}}\]
\[\kf{= 1\cdot(-1)(0-8) + (-1)\cdot 1\cdot(2-12) + 3\cdot(-1)(4-0)}\]
\[\kf{= 8 + 10 - 12 = 6,}\]
which agrees with the value found by Sarrus’ rule in the previous section.
\end{misol}

\bigskip
\subsection{Methods for evaluating determinants of order \texorpdfstring{$n$}{n}}
\textbf{Repeated Laplace expansion.}
Expanding a determinant of order $n$ along a row (column) reduces the problem to
determinants of order $n-1$. Repeating the procedure, we eventually arrive at
determinants of order two.
\vspace{3pt}
\begin{misol}
Consider the fourth-order determinant
\[\kf{D = \begin{vmatrix}
1 & 0 & 2 & -1 \\
3 & 0 & 0 & 5 \\
2 & 1 & 4 & -3 \\
1 & 0 & 0 & 0
\end{vmatrix}.}\]
The fourth row contains three zeros, so we expand along it:
\[\kf{D = 1\cdot(-1)^{4+1}\begin{vmatrix}0&2&-1\\0&0&5\\1&4&-3\end{vmatrix}
= -\begin{vmatrix}0&2&-1\\0&0&5\\1&4&-3\end{vmatrix}.}\]
The resulting $3\times3$ determinant is expanded along its first column:
\[\kf{\begin{vmatrix}0&2&-1\\0&0&5\\1&4&-3\end{vmatrix}
= 1\cdot(-1)^{3+1}\begin{vmatrix}2&-1\\0&5\end{vmatrix}
= 1\cdot(10-0) = 10.}\]
Hence $D = -10$.
\end{misol}
\textbf{Reduction by elementary operations.}
The most efficient approach is to use the properties of determinants to create
zeros: adding suitable multiples of one row (column) to the others, we leave a
single nonzero entry in some column (row) and then expand along it. The order of
the determinant drops by one, and we repeat the process until a $2\times2$
determinant remains.
\begin{misol}
Evaluate
\[\kf{D = \begin{vmatrix}
1 & 2 & -1 & 3 \\
2 & 5 & 0 & -1 \\
-1 & 0 & 3 & 2 \\
3 & 1 & -2 & 4
\end{vmatrix}.}\]
We clear the first column by the row operations $S_2 - 2S_1$, $S_3 + S_1$,
$S_4 - 3S_1$:
\[\kf{D = \begin{vmatrix}
1 & 2 & -1 & 3 \\
0 & 1 & 2 & -7 \\
0 & 2 & 2 & 5 \\
0 & -5 & 1 & -5
\end{vmatrix}.}\]
Expanding along the first column gives $D = 1\cdot(-1)^{1+1}\cdot D_3$, where
\[\kf{D_3 = \begin{vmatrix}
1 & 2 & -7 \\
2 & 2 & 5 \\
-5 & 1 & -5
\end{vmatrix}.}\]
Clearing the first column of $D_3$ by $S_2 - 2S_1$ and $S_3 + 5S_1$, we get
\[\kf{D_3 = \begin{vmatrix}
1 & 2 & -7 \\
0 & -2 & 19 \\
0 & 11 & -40
\end{vmatrix}
= 1\cdot\begin{vmatrix}-2 & 19 \\ 11 & -40\end{vmatrix}
= (-2)(-40) - 19\cdot 11 = 80 - 209 = -129.}\]
Therefore $D = -129$.
\end{misol}
\begin{misol}
Evaluate the fifth-order determinant
\[\kf{D = \begin{vmatrix}
1 & 0 & 0 & 2 & 0 \\
3 & 1 & 0 & -1 & 0 \\
0 & 2 & 1 & 4 & 0 \\
-1 & 0 & 3 & 0 & 0 \\
2 & 1 & 0 & 3 & 5
\end{vmatrix}.}\]
The fifth column has only one nonzero entry, and expanding along it we obtain
\[\kf{D = 5\cdot(-1)^{5+5}\begin{vmatrix}
1 & 0 & 0 & 2 \\
3 & 1 & 0 & -1 \\
0 & 2 & 1 & 4 \\
-1 & 0 & 3 & 0
\end{vmatrix}
= 5\cdot D_4.}\]
Next we expand $D_4$ along its first row:
\[\kf{D_4 = 1\cdot(-1)^{1+1}\begin{vmatrix}1&0&-1\\2&1&4\\0&3&0\end{vmatrix}
+ 2\cdot(-1)^{1+4}\begin{vmatrix}3&1&0\\0&2&1\\-1&0&3\end{vmatrix}.}\]
The first of these third-order determinants is expanded along its third row,
\[\kf{\begin{vmatrix}1&0&-1\\2&1&4\\0&3&0\end{vmatrix}
= 3\cdot(-1)^{3+2}\begin{vmatrix}1&-1\\2&4\end{vmatrix}
= -3\cdot(4+2) = -18,}\]
and the second along its first row,
\[\kf{\begin{vmatrix}3&1&0\\0&2&1\\-1&0&3\end{vmatrix}
= 3(6-0) - 1(0+1) + 0 = 18 - 1 = 17.}\]
Thus $D_4 = -18 - 2\cdot 17 = -52$ and $D = 5\cdot(-52) = -260$.
\end{misol}
\bigskip
\subsection{The determinant of a triangular matrix}
\vspace{3pt}
\begin{teorema}
The determinant of a triangular (upper or lower) matrix equals the product of its
\textbf{main diagonal} entries:
\[\kf{\begin{vmatrix}
a_{11} & a_{12} & \cdots & a_{1n} \\
0 & a_{22} & \cdots & a_{2n} \\
\vdots & \vdots & \ddots & \vdots \\
0 & 0 & \cdots & a_{nn}
\end{vmatrix}
= a_{11}\cdot a_{22}\cdot \ldots \cdot a_{nn}.}\]
\end{teorema}
To see this, expand repeatedly along the first column, where only the diagonal
entry is nonzero. The theorem underlies the method of elementary operations: we
bring the determinant to triangular form and multiply the diagonal entries.
\vspace{3pt}
\begin{misol}
\[\kf{\begin{vmatrix}
3 & 5 & -1 & 2 \\
0 & -2 & 4 & 1 \\
0 & 0 & 7 & -3 \\
0 & 0 & 0 & 4
\end{vmatrix}
= 3\cdot(-2)\cdot 7\cdot 4 = -168.}\]
\end{misol}
\bigskip
% =========================================================
\subsection{Minors and cofactors in information security}
\textbf{The inverse matrix and decryption.}
Decryption in the Hill cipher requires the inverse of the key matrix. The formula
for the inverse is built directly from \textbf{cofactors}:
\[\kf{A^{-1} = \frac{1}{\det A}\cdot \widetilde{A}^{T},}\]
where $\widetilde{A} = (A_{ij})$ is the matrix of cofactors. Each $A_{ij}$ is computed
from the minor $M_{ij}$, so minors and cofactors lie at the heart of the mechanism
by which a cryptosystem “reverses” its key.

\textbf{Minors in coding theory.}
For error-correcting codes (Hamming codes, Reed–Solomon codes and others) to work,
the rows of the generator matrix $G$ and of the parity-check matrix $H$ must be
linearly independent. This is verified by computing various \textbf{minors} of the
matrix: if at least one minor of order $k$ is nonzero, the rank of the matrix is
at least $k$.

\textbf{Secret sharing.}
In \textbf{Shamir’s secret sharing scheme} a secret key is distributed among $n$
participants in such a way that any $k$ of them can recover it, while fewer
cannot. Recovery is based on Lagrange interpolation, which in turn is tied to
Vandermonde determinants; the scheme works precisely because these determinants
are nonzero for distinct interpolation nodes.

\textbf{Digital signatures and authentication.}
Some digital signature schemes, in particular matrix-based ones, verify a signature
by computing determinants and minors of certain matrices; the signature is
accepted if the determinant takes the expected value.

\subsection{Computing with Python, Maple and Mathcad}
We return to the matrix $A$ from the first examples of this section and compute
the minor $M_{12}$, the cofactor $A_{12}$ and the whole matrix of cofactors, and
then the fourth-order determinant evaluated above by elementary operations.
Note that SymPy numbers rows and columns from $0$, so $M_{12}$ is
\texttt{A.minor(0, 1)}.
\begin{Verbatim}[fontsize=\footnotesize,frame=single,label={\small Python}]
from sympy import Matrix
A = Matrix([[2, 1, 3], [0, -1, 2], [4, 3, 1]])
print(A.minor(0, 1), A.cofactor(0, 1))
print(A.cofactor_matrix())
D = Matrix([[1, 2, -1, 3], [2, 5, 0, -1],
            [-1, 0, 3, 2], [3, 1, -2, 4]])
print(D.det())
\end{Verbatim}
\begin{Verbatim}[fontsize=\footnotesize,frame=single,label={\small Output}]
-8 8
Matrix([[-7, 8, 4], [8, -10, -2], [5, -4, -2]])
-129
\end{Verbatim}
In Maple, \texttt{Minor(A, i, j)} returns the minor $M_{ij}$ (numbering starts
at $1$); the cofactor is obtained by attaching the sign:
\begin{Verbatim}[fontsize=\footnotesize,frame=single,label={\small Maple}]
with(LinearAlgebra):
A := Matrix([[2, 1, 3], [0, -1, 2], [4, 3, 1]]):
Minor(A, 1, 2);
                              -8
(-1)^(1+2)*Minor(A, 1, 2);
                               8
D4 := Matrix([[1, 2, -1, 3], [2, 5, 0, -1],
              [-1, 0, 3, 2], [3, 1, -2, 4]]):
Determinant(D4);
                             -129
\end{Verbatim}
Mathcad Prime has no separate minor function, but a submatrix is easy to form:
\texttt{submatrix(A, ir, jr, ic, jc)} extracts rows \texttt{ir..jr} and columns
\texttt{ic..jc}, and \texttt{augment} glues columns together. With the default
\texttt{ORIGIN = 0}, the minor $M_{12}$ is the determinant of rows 1–2 and
columns 0 and 2:
\begin{Verbatim}[fontsize=\footnotesize,frame=single,label={\small Mathcad}]
     | 2   1  3 |
A := | 0  -1  2 |
     | 4   3  1 |
B := submatrix(A, 1, 2, 0, 2)
M12 := |augment(B<0>, B<2>)|      M12 = -8
A12 := (-1)^(1+2) * M12           A12 = 8

     |  1  2  -1   3 |
D := |  2  5   0  -1 |            |D| = -129
     | -1  0   3   2 |
     |  3  1  -2   4 |
\end{Verbatim}
Here \texttt{B<0>} denotes the column operator (column 0 of $B$).

\subsection{Exercises on minors and cofactors}
\bigskip
\emph{1. For the given matrix, find the minors $M_{ij}$ and the cofactors $A_{ij}$ of the indicated entries.}
\bigskip
\begin{enumerate}[label=\arabic*.,  leftmargin=*]
\item $A=\begin{pmatrix}3&-1&2\\0&4&-5\\1&2&3\end{pmatrix}$\quad $M_{11}, A_{11}, M_{23}, A_{23}$.
\item $A=\begin{pmatrix}-2&1&3\\4&0&-1\\5&2&-4\end{pmatrix}$\quad $M_{12}, A_{12}, M_{31}, A_{31}$.
\item $A=\begin{pmatrix}1&5&-3\\-2&0&4\\3&1&2\end{pmatrix}$\quad $M_{22}, A_{22}, M_{13}, A_{13}$.
\item $A=\begin{pmatrix}4&-3&0\\2&1&-5\\-1&6&3\end{pmatrix}$\quad $M_{32}, A_{32}, M_{11}, A_{11}$.
\item $A=\begin{pmatrix}0&2&-1\\3&-4&5\\1&0&-2\end{pmatrix}$\quad $M_{21}, A_{21}, M_{33}, A_{33}$.
\item $A=\begin{pmatrix}-5&1&4\\0&3&-2\\2&-1&0\end{pmatrix}$\quad $M_{13}, A_{13}, M_{22}, A_{22}$.
\item $A=\begin{pmatrix}2&0&-3\\-1&5&1\\4&-2&3\end{pmatrix}$\quad $M_{31}, A_{31}, M_{12}, A_{12}$.
\item $A=\begin{pmatrix}1&-4&2\\3&0&-5\\-2&1&4\end{pmatrix}$\quad $M_{23}, A_{23}, M_{11}, A_{11}$.
\item $A=\begin{pmatrix}-3&2&0\\5&-1&4\\0&3&-2\end{pmatrix}$\quad $M_{32}, A_{32}, M_{21}, A_{21}$.
\item $A=\begin{pmatrix}6&-2&1\\0&4&-3\\3&1&5\end{pmatrix}$\quad $M_{11}, A_{11}, M_{33}, A_{33}$.
\item $A=\begin{pmatrix}-1&3&5\\2&0&-4\\4&-5&1\end{pmatrix}$\quad $M_{12}, A_{12}, M_{31}, A_{31}$.
\item $A=\begin{pmatrix}3&-5&0\\-1&2&4\\5&0&-3\end{pmatrix}$\quad $M_{22}, A_{22}, M_{13}, A_{13}$.
\item $A=\begin{pmatrix}0&4&-2\\1&-3&5\\-4&2&0\end{pmatrix}$\quad $M_{21}, A_{21}, M_{32}, A_{32}$.
\item $A=\begin{pmatrix}2&1&-4\\-3&5&0\\0&-2&3\end{pmatrix}$\quad $M_{13}, A_{13}, M_{23}, A_{23}$.
\item $A=\begin{pmatrix}-4&0&3\\2&-1&5\\1&4&-2\end{pmatrix}$\quad $M_{11}, A_{11}, M_{32}, A_{32}$.
\item $A=\begin{pmatrix}5&-3&1\\0&2&-4\\-2&1&0\end{pmatrix}$\quad $M_{31}, A_{31}, M_{12}, A_{12}$.
\item $A=\begin{pmatrix}1&0&-5\\-4&3&2\\3&-1&4\end{pmatrix}$\quad $M_{22}, A_{22}, M_{33}, A_{33}$.
\item $A=\begin{pmatrix}-2&4&0\\1&-5&3\\0&2&-1\end{pmatrix}$\quad $M_{13}, A_{13}, M_{21}, A_{21}$.
\item $A=\begin{pmatrix}3&1&-2\\-5&0&4\\2&-3&1\end{pmatrix}$\quad $M_{23}, A_{23}, M_{11}, A_{11}$.
\item $A=\begin{pmatrix}0&-2&3\\4&1&-5\\-1&5&2\end{pmatrix}$\quad $M_{32}, A_{32}, M_{12}, A_{12}$.
\item $A=\begin{pmatrix}-1&5&0\\3&-2&4\\2&0&-3\end{pmatrix}$\quad $M_{11}, A_{11}, M_{23}, A_{23}$.
\item $A=\begin{pmatrix}4&0&-1\\-2&3&5\\1&-4&2\end{pmatrix}$\quad $M_{21}, A_{21}, M_{33}, A_{33}$.
\item $A=\begin{pmatrix}-3&1&4\\0&-5&2\\5&3&0\end{pmatrix}$\quad $M_{31}, A_{31}, M_{22}, A_{22}$.
\item $A=\begin{pmatrix}2&-4&1\\5&0&-3\\-1&2&4\end{pmatrix}$\quad $M_{12}, A_{12}, M_{33}, A_{33}$.
\item $A=\begin{pmatrix}1&3&-5\\-4&0&2\\0&-1&3\end{pmatrix}$\quad $M_{13}, A_{13}, M_{21}, A_{21}$.
\item $A=\begin{pmatrix}-5&2&0\\1&-3&4\\3&0&-1\end{pmatrix}$\quad $M_{22}, A_{22}, M_{31}, A_{31}$.
\item $A=\begin{pmatrix}0&-1&3\\2&4&0\\-3&5&-2\end{pmatrix}$\quad $M_{32}, A_{32}, M_{11}, A_{11}$.
\item $A=\begin{pmatrix}3&0&-4\\-1&2&5\\4&-3&1\end{pmatrix}$\quad $M_{23}, A_{23}, M_{12}, A_{12}$.
\item $A=\begin{pmatrix}-2&5&1\\0&-4&3\\1&2&0\end{pmatrix}$\quad $M_{11}, A_{11}, M_{13}, A_{13}$.
\item $A=\begin{pmatrix}4&-1&0\\-3&0&2\\2&5&-4\end{pmatrix}$\quad $M_{21}, A_{21}, M_{32}, A_{32}$.
\end{enumerate}
\emph{2. Evaluate the determinant by expanding along the indicated row or column.}
\bigskip
\begin{multicols}{2}\raggedright
\begin{enumerate}[label=\arabic*.,  leftmargin=*]
\item $\begin{vmatrix}1&0&3\\2&-1&4\\0&5&-2\end{vmatrix}$, \quad row 1.
\item $\begin{vmatrix}-2&3&0\\1&0&-4\\5&2&1\end{vmatrix}$, \quad column 2.
\item $\begin{vmatrix}4&-1&2\\0&3&0\\-2&5&1\end{vmatrix}$, \quad row 2.
\item $\begin{vmatrix}0&2&-3\\4&0&1\\-1&5&0\end{vmatrix}$, \quad column 1.
\item $\begin{vmatrix}3&0&-1\\0&4&2\\5&-3&0\end{vmatrix}$, \quad column 3.
\item $\begin{vmatrix}-1&2&5\\3&0&0\\4&-2&1\end{vmatrix}$, \quad row 2.
\item $\begin{vmatrix}2&-3&0\\0&1&4\\-5&0&3\end{vmatrix}$, \quad row 3.
\item $\begin{vmatrix}0&5&-1\\2&0&3\\-4&1&0\end{vmatrix}$, \quad row 1.
\item $\begin{vmatrix}1&-2&0\\0&3&-5\\4&0&2\end{vmatrix}$, \quad column 3.
\item $\begin{vmatrix}-3&0&4\\2&1&0\\0&-5&3\end{vmatrix}$, \quad row 2.
\item $\begin{vmatrix}5&0&0\\-1&2&-3\\3&4&1\end{vmatrix}$, \quad row 1.
\item $\begin{vmatrix}0&-4&1\\3&2&0\\-1&0&5\end{vmatrix}$, \quad column 2.
\item $\begin{vmatrix}2&1&0\\-3&0&4\\0&5&-2\end{vmatrix}$, \quad column 1.
\item $\begin{vmatrix}0&0&3\\-2&4&1\\5&-1&0\end{vmatrix}$, \quad row 1.
\item $\begin{vmatrix}-1&3&0\\0&2&-4\\4&0&5\end{vmatrix}$, \quad column 3.
\item $\begin{vmatrix}3&-2&1\\0&0&-5\\2&4&0\end{vmatrix}$, \quad row 2.
\item $\begin{vmatrix}0&1&-3\\5&0&2\\-2&4&0\end{vmatrix}$, \quad column 1.
\item $\begin{vmatrix}-4&0&2\\1&3&0\\0&-5&4\end{vmatrix}$, \quad row 3.
\item $\begin{vmatrix}1&0&-2\\0&-3&5\\3&0&4\end{vmatrix}$, \quad column 2.
\item $\begin{vmatrix}2&5&0\\0&-1&3\\-4&0&2\end{vmatrix}$, \quad column 3.
\item $\begin{vmatrix}0&-2&4\\3&1&0\\-5&0&2\end{vmatrix}$, \quad row 2.
\item $\begin{vmatrix}-1&0&5\\4&2&0\\0&-3&1\end{vmatrix}$, \quad row 1.
\item $\begin{vmatrix}3&0&-4\\0&5&1\\2&0&3\end{vmatrix}$, \quad column 2.
\item $\begin{vmatrix}0&4&-1\\-2&0&3\\5&1&0\end{vmatrix}$, \quad row 3.
\item $\begin{vmatrix}1&-5&0\\0&2&3\\-3&0&4\end{vmatrix}$, \quad column 1.
\item $\begin{vmatrix}-2&0&3\\4&-1&0\\0&5&-2\end{vmatrix}$, \quad column 3.
\item $\begin{vmatrix}5&1&0\\0&-3&2\\-4&0&4\end{vmatrix}$, \quad row 3.
\item $\begin{vmatrix}0&2&-5\\1&0&4\\-3&3&0\end{vmatrix}$, \quad column 2.
\item $\begin{vmatrix}-1&0&0\\3&-2&5\\2&4&-3\end{vmatrix}$, \quad row 1.
\item $\begin{vmatrix}4&-3&0\\0&1&-2\\2&0&5\end{vmatrix}$, \quad row 2.
\end{enumerate}
\end{multicols}
\emph{3. Evaluate the fourth-order determinants.}
\bigskip
\begin{multicols}{3}\raggedright
\begin{enumerate}[label=\arabic*.,  leftmargin=*]
\item $\begin{vmatrix}1&0&2&-1\\3&0&0&5\\2&1&4&-3\\1&0&0&0\end{vmatrix}$
\item $\begin{vmatrix}2&-1&0&3\\0&4&1&0\\-1&0&3&2\\0&5&0&-4\end{vmatrix}$
\item $\begin{vmatrix}1&2&-1&0\\0&3&0&-2\\4&0&1&5\\-1&0&2&0\end{vmatrix}$
\item $\begin{vmatrix}-2&0&1&0\\3&-1&0&4\\0&2&-3&0\\1&0&0&5\end{vmatrix}$
\item $\begin{vmatrix}0&1&-2&3\\1&0&4&0\\-3&2&0&1\\0&0&5&-1\end{vmatrix}$
\item $\begin{vmatrix}3&0&0&-2\\-1&2&0&4\\0&5&1&0\\4&0&-3&0\end{vmatrix}$
\item $\begin{vmatrix}1&-3&0&0\\0&2&-1&0\\4&0&3&-2\\0&1&0&5\end{vmatrix}$
\item $\begin{vmatrix}-1&0&4&0\\2&3&0&-1\\0&0&5&2\\3&-2&0&0\end{vmatrix}$
\item $\begin{vmatrix}2&1&0&-3\\0&-2&3&0\\1&0&0&4\\0&4&-1&0\end{vmatrix}$
\item $\begin{vmatrix}0&0&2&1\\-3&1&0&0\\0&4&-1&3\\5&0&0&-2\end{vmatrix}$
\item $\begin{vmatrix}1&2&3&4\\0&1&2&3\\0&0&1&2\\0&0&0&1\end{vmatrix}$
\item $\begin{vmatrix}-1&0&0&3\\2&-4&0&0\\0&1&5&0\\3&0&-2&1\end{vmatrix}$
\item $\begin{vmatrix}4&0&-1&0\\0&3&0&2\\-2&0&5&0\\0&1&0&-3\end{vmatrix}$
\item $\begin{vmatrix}1&-1&2&0\\3&0&0&-4\\0&2&1&3\\-2&0&0&5\end{vmatrix}$
\item $\begin{vmatrix}0&3&0&-1\\2&0&-4&0\\0&1&0&3\\-5&0&2&0\end{vmatrix}$
\item $\begin{vmatrix}2&0&0&1\\0&-3&0&0\\4&0&1&-2\\0&2&0&3\end{vmatrix}$
\item $\begin{vmatrix}1&0&-2&3\\0&2&0&0\\-1&0&4&-1\\3&0&0&2\end{vmatrix}$
\item $\begin{vmatrix}-3&1&0&0\\0&0&2&-4\\5&-1&0&0\\0&0&3&1\end{vmatrix}$
\item $\begin{vmatrix}0&2&-1&0\\1&0&3&-2\\0&-4&0&5\\3&0&1&0\end{vmatrix}$
\item $\begin{vmatrix}1&0&0&0\\-2&3&0&0\\4&-1&2&0\\0&5&-3&1\end{vmatrix}$
\item $\begin{vmatrix}2&-1&0&4\\0&3&-2&0\\1&0&0&-3\\0&0&5&1\end{vmatrix}$
\item $\begin{vmatrix}0&0&1&-2\\0&4&0&3\\-1&0&0&5\\2&-3&0&0\end{vmatrix}$
\item $\begin{vmatrix}5&0&-1&0\\0&2&0&-3\\3&0&4&0\\0&-1&0&2\end{vmatrix}$
\item $\begin{vmatrix}-2&3&0&0\\1&0&-4&0\\0&2&0&5\\0&0&1&-3\end{vmatrix}$
\item $\begin{vmatrix}1&0&3&0\\0&-2&0&4\\2&0&-1&0\\0&3&0&5\end{vmatrix}$
\item $\begin{vmatrix}0&1&0&-3\\4&0&2&0\\0&-5&0&1\\3&0&-1&0\end{vmatrix}$
\item $\begin{vmatrix}3&0&0&-1\\0&2&4&0\\0&-1&3&0\\5&0&0&2\end{vmatrix}$
\item $\begin{vmatrix}-1&2&0&0\\0&0&3&-5\\4&-1&0&0\\0&0&2&1\end{vmatrix}$
\item $\begin{vmatrix}2&0&-3&0\\0&1&0&4\\-1&0&2&0\\0&-3&0&5\end{vmatrix}$
\item $\begin{vmatrix}0&-2&0&1\\3&0&5&0\\0&4&0&-2\\1&0&-3&0\end{vmatrix}$
\end{enumerate}
\end{multicols}
\bigskip

\emph{4. Using cofactors, compute the sum $\sum\limits_{j=1}^{3}a_{ij}A_{kj}$ in the cases $i=k$ and $i\neq k$.}
\begin{misol}
For $A=\begin{pmatrix}2&-1&3\\4&0&-5\\1&3&2\end{pmatrix}$ compute
1) $\sum\limits_{j=1}^{3}a_{1j}A_{1j}$ (the case $i=1,\,k=1$);
2) $\sum\limits_{j=1}^{3}a_{1j}A_{2j}$ (the case $i=1,\,k=2$).

\textbf{Solution.} The cofactors of the first row are
\[\kf{A_{11}=\begin{vmatrix}0&-5\\3&2\end{vmatrix}=15,\quad
A_{12}=-\begin{vmatrix}4&-5\\1&2\end{vmatrix}=-13,\quad
A_{13}=\begin{vmatrix}4&0\\1&3\end{vmatrix}=12.}\]
In the first case ($i=1,\,k=1$) we obtain
\[\kf{\sum\limits_{j=1}^{3} a_{1j}A_{1j}=a_{11}A_{11} + a_{12}A_{12} + a_{13}A_{13} = 2\cdot15+(-1)(-13)+3\cdot12=}\]
\[\kf{=30+13+36=79=\det A.}\]
For the second case we need the cofactors of the second row:
\[\kf{A_{21}=-\begin{vmatrix}-1&3\\3&2\end{vmatrix}=11,\quad
A_{22}=\begin{vmatrix}2&3\\1&2\end{vmatrix}=1,\quad
A_{23}=-\begin{vmatrix}2&-1\\1&3\end{vmatrix}=-7.}\]
Then ($i=1,\,k=2$)
\[\kf{\sum\limits_{j=1}^{3} a_{1j}A_{2j}=a_{11}A_{21} + a_{12}A_{22} + a_{13}A_{23} = 2\cdot11+(-1)\cdot1+3\cdot(-7)=}\]
\[\kf{=22-1-21=0.}\]
This illustrates the general rule: the entries of a row multiplied by their own
cofactors ($i=k$) add up to $\det A$, whereas multiplied by the cofactors of
another row ($i\neq k$) they always add up to $0$.
\end{misol}
\bigskip
\begin{enumerate}[label=\arabic*.,  leftmargin=*]
\item $A=\begin{pmatrix}2&-1&3\\4&0&-5\\1&3&2\end{pmatrix}$; \quad $i=1, k=1$;\quad $i=1, k=2$.
\item $A=\begin{pmatrix}-3&2&1\\0&4&-2\\5&-1&3\end{pmatrix}$; \quad $i=2, k=2$;\quad $i=2, k=3$.
\item $A=\begin{pmatrix}1&0&-4\\3&-2&5\\-1&4&0\end{pmatrix}$; \quad $i=3, k=3$;\quad $i=3, k=1$.
\item $A=\begin{pmatrix}0&5&-1\\2&-3&4\\-4&1&2\end{pmatrix}$; \quad $i=1, k=1$;\quad $i=1, k=3$.
\item $A=\begin{pmatrix}-2&1&0\\3&0&-5\\1&4&3\end{pmatrix}$; \quad $i=2, k=2$;\quad $i=2, k=1$.
\item $A=\begin{pmatrix}4&-3&2\\0&1&-4\\-2&5&0\end{pmatrix}$; \quad $i=3, k=3$;\quad $i=3, k=2$.
\item $A=\begin{pmatrix}1&-2&5\\-3&4&0\\2&0&-1\end{pmatrix}$; \quad $i=1, k=1$;\quad $i=1, k=3$.
\item $A=\begin{pmatrix}-5&0&3\\1&-4&2\\0&3&-2\end{pmatrix}$; \quad $i=2, k=2$;\quad $i=2, k=3$.
\item $A=\begin{pmatrix}3&1&-2\\-1&5&0\\4&-3&1\end{pmatrix}$; \quad $i=3, k=3$;\quad $i=3, k=1$.
\item $A=\begin{pmatrix}0&-4&2\\5&1&-3\\-2&0&4\end{pmatrix}$; \quad $i=1, k=1$;\quad $i=1, k=2$.
\item $A=\begin{pmatrix}-1&3&0\\2&-5&4\\0&1&-3\end{pmatrix}$; \quad $i=2, k=2$;\quad $i=2, k=1$.
\item $A=\begin{pmatrix}5&0&-3\\-2&4&1\\3&-1&0\end{pmatrix}$; \quad $i=3, k=3$;\quad $i=3, k=2$.
\item $A=\begin{pmatrix}2&-1&4\\0&3&-2\\-5&2&1\end{pmatrix}$; \quad $i=1, k=1$;\quad $i=1, k=3$.
\item $A=\begin{pmatrix}-4&5&0\\1&0&-3\\3&-2&4\end{pmatrix}$; \quad $i=2, k=2$;\quad $i=2, k=3$.
\item $A=\begin{pmatrix}0&2&-5\\-3&1&4\\4&0&-1\end{pmatrix}$; \quad $i=3, k=3$;\quad $i=3, k=1$.
\item $A=\begin{pmatrix}1&-3&2\\4&0&-1\\-2&5&3\end{pmatrix}$; \quad $i=1, k=1$;\quad $i=1, k=2$.
\item $A=\begin{pmatrix}-2&0&5\\3&-4&1\\0&2&-3\end{pmatrix}$; \quad $i=2, k=2$;\quad $i=2, k=1$.
\item $A=\begin{pmatrix}3&4&0\\-1&-2&5\\2&0&-4\end{pmatrix}$; \quad $i=3, k=3$;\quad $i=3, k=2$.
\item $A=\begin{pmatrix}0&-5&1\\2&3&-4\\-1&0&2\end{pmatrix}$; \quad $i=1, k=1$;\quad $i=1, k=3$.
\item $A=\begin{pmatrix}-3&1&0\\0&-2&3\\5&4&-1\end{pmatrix}$; \quad $i=2, k=2$;\quad $i=2, k=3$.
\item $A=\begin{pmatrix}4&0&-2\\-1&3&5\\0&-4&1\end{pmatrix}$; \quad $i=3, k=3$;\quad $i=3, k=1$.
\item $A=\begin{pmatrix}1&2&-3\\-4&0&5\\3&-1&0\end{pmatrix}$; \quad $i=1, k=1$;\quad $i=1, k=2$.
\item $A=\begin{pmatrix}-5&3&0\\2&-1&4\\0&5&-2\end{pmatrix}$; \quad $i=2, k=2$;\quad $i=2, k=1$.
\item $A=\begin{pmatrix}0&-3&4\\1&5&0\\-2&0&3\end{pmatrix}$; \quad $i=3, k=3$;\quad $i=3, k=2$.
\item $A=\begin{pmatrix}2&0&1\\-3&-4&0\\5&2&-5\end{pmatrix}$; \quad $i=1, k=1$;\quad $i=1, k=3$.
\item $A=\begin{pmatrix}-1&4&-2\\0&0&3\\3&-5&1\end{pmatrix}$; \quad $i=2, k=2$;\quad $i=2, k=3$.
\item $A=\begin{pmatrix}3&-2&0\\-4&1&5\\0&3&-1\end{pmatrix}$; \quad $i=3, k=3$;\quad $i=3, k=1$.
\item $A=\begin{pmatrix}0&1&-4\\5&-3&2\\-1&0&3\end{pmatrix}$; \quad $i=1, k=1$;\quad $i=1, k=2$.
\item $A=\begin{pmatrix}-2&0&3\\1&5&-4\\4&-1&0\end{pmatrix}$; \quad $i=2, k=2$;\quad $i=2, k=1$.
\item $A=\begin{pmatrix}5&-4&1\\0&2&-3\\-3&0&4\end{pmatrix}$; \quad $i=3, k=3$;\quad $i=3, k=2$.
\end{enumerate}